\IfFileExists{stacks-project.cls}{%
\documentclass{stacks-project}
}{%
\documentclass{amsart}
}

% For dealing with references we use the comment environment
\usepackage{verbatim}
\newenvironment{reference}{\comment}{\endcomment}
%\newenvironment{reference}{}{}
\newenvironment{slogan}{\comment}{\endcomment}
\newenvironment{history}{\comment}{\endcomment}

% For commutative diagrams we use Xy-pic
\usepackage[all]{xy}

% We use 2cell for 2-commutative diagrams.
\xyoption{2cell}
\UseAllTwocells

% We use multicol for the list of chapters between chapters
\usepackage{multicol}

% This is generally recommended for better output
\usepackage{lmodern}
\usepackage[T1]{fontenc}

% For cross-file-references
\usepackage{xr-hyper}

% Package for hypertext links:
\usepackage{hyperref}

% For any local file, say "hello.tex" you want to link to please
% use \externaldocument[hello-]{hello}
\externaldocument[introduction-]{introduction}
\externaldocument[conventions-]{conventions}
\externaldocument[sets-]{sets}
\externaldocument[categories-]{categories}
\externaldocument[topology-]{topology}
\externaldocument[sheaves-]{sheaves}
\externaldocument[sites-]{sites}
\externaldocument[stacks-]{stacks}
\externaldocument[fields-]{fields}
\externaldocument[algebra-]{algebra}
\externaldocument[brauer-]{brauer}
\externaldocument[homology-]{homology}
\externaldocument[derived-]{derived}
\externaldocument[simplicial-]{simplicial}
\externaldocument[more-algebra-]{more-algebra}
\externaldocument[smoothing-]{smoothing}
\externaldocument[modules-]{modules}
\externaldocument[sites-modules-]{sites-modules}
\externaldocument[injectives-]{injectives}
\externaldocument[cohomology-]{cohomology}
\externaldocument[sites-cohomology-]{sites-cohomology}
\externaldocument[dga-]{dga}
\externaldocument[dpa-]{dpa}
\externaldocument[sdga-]{sdga}
\externaldocument[hypercovering-]{hypercovering}
\externaldocument[schemes-]{schemes}
\externaldocument[constructions-]{constructions}
\externaldocument[properties-]{properties}
\externaldocument[morphisms-]{morphisms}
\externaldocument[coherent-]{coherent}
\externaldocument[divisors-]{divisors}
\externaldocument[limits-]{limits}
\externaldocument[varieties-]{varieties}
\externaldocument[topologies-]{topologies}
\externaldocument[descent-]{descent}
\externaldocument[perfect-]{perfect}
\externaldocument[more-morphisms-]{more-morphisms}
\externaldocument[flat-]{flat}
\externaldocument[groupoids-]{groupoids}
\externaldocument[more-groupoids-]{more-groupoids}
\externaldocument[etale-]{etale}
\externaldocument[chow-]{chow}
\externaldocument[intersection-]{intersection}
\externaldocument[pic-]{pic}
\externaldocument[weil-]{weil}
\externaldocument[adequate-]{adequate}
\externaldocument[dualizing-]{dualizing}
\externaldocument[duality-]{duality}
\externaldocument[discriminant-]{discriminant}
\externaldocument[derham-]{derham}
\externaldocument[local-cohomology-]{local-cohomology}
\externaldocument[algebraization-]{algebraization}
\externaldocument[curves-]{curves}
\externaldocument[resolve-]{resolve}
\externaldocument[models-]{models}
\externaldocument[functors-]{functors}
\externaldocument[equiv-]{equiv}
\externaldocument[pione-]{pione}
\externaldocument[etale-cohomology-]{etale-cohomology}
\externaldocument[proetale-]{proetale}
\externaldocument[relative-cycles-]{relative-cycles}
\externaldocument[more-etale-]{more-etale}
\externaldocument[trace-]{trace}
\externaldocument[crystalline-]{crystalline}
\externaldocument[spaces-]{spaces}
\externaldocument[spaces-properties-]{spaces-properties}
\externaldocument[spaces-morphisms-]{spaces-morphisms}
\externaldocument[decent-spaces-]{decent-spaces}
\externaldocument[spaces-cohomology-]{spaces-cohomology}
\externaldocument[spaces-limits-]{spaces-limits}
\externaldocument[spaces-divisors-]{spaces-divisors}
\externaldocument[spaces-over-fields-]{spaces-over-fields}
\externaldocument[spaces-topologies-]{spaces-topologies}
\externaldocument[spaces-descent-]{spaces-descent}
\externaldocument[spaces-perfect-]{spaces-perfect}
\externaldocument[spaces-more-morphisms-]{spaces-more-morphisms}
\externaldocument[spaces-flat-]{spaces-flat}
\externaldocument[spaces-groupoids-]{spaces-groupoids}
\externaldocument[spaces-more-groupoids-]{spaces-more-groupoids}
\externaldocument[bootstrap-]{bootstrap}
\externaldocument[spaces-pushouts-]{spaces-pushouts}
\externaldocument[spaces-chow-]{spaces-chow}
\externaldocument[groupoids-quotients-]{groupoids-quotients}
\externaldocument[spaces-more-cohomology-]{spaces-more-cohomology}
\externaldocument[spaces-simplicial-]{spaces-simplicial}
\externaldocument[spaces-duality-]{spaces-duality}
\externaldocument[formal-spaces-]{formal-spaces}
\externaldocument[restricted-]{restricted}
\externaldocument[spaces-resolve-]{spaces-resolve}
\externaldocument[formal-defos-]{formal-defos}
\externaldocument[defos-]{defos}
\externaldocument[cotangent-]{cotangent}
\externaldocument[examples-defos-]{examples-defos}
\externaldocument[algebraic-]{algebraic}
\externaldocument[examples-stacks-]{examples-stacks}
\externaldocument[stacks-sheaves-]{stacks-sheaves}
\externaldocument[criteria-]{criteria}
\externaldocument[artin-]{artin}
\externaldocument[quot-]{quot}
\externaldocument[stacks-properties-]{stacks-properties}
\externaldocument[stacks-morphisms-]{stacks-morphisms}
\externaldocument[stacks-limits-]{stacks-limits}
\externaldocument[stacks-cohomology-]{stacks-cohomology}
\externaldocument[stacks-perfect-]{stacks-perfect}
\externaldocument[stacks-introduction-]{stacks-introduction}
\externaldocument[stacks-more-morphisms-]{stacks-more-morphisms}
\externaldocument[stacks-geometry-]{stacks-geometry}
\externaldocument[moduli-]{moduli}
\externaldocument[moduli-curves-]{moduli-curves}
\externaldocument[examples-]{examples}
\externaldocument[exercises-]{exercises}
\externaldocument[guide-]{guide}
\externaldocument[desirables-]{desirables}
\externaldocument[coding-]{coding}
\externaldocument[obsolete-]{obsolete}
\externaldocument[fdl-]{fdl}
\externaldocument[index-]{index}

% Theorem environments.
%
\theoremstyle{plain}
\newtheorem{theorem}[subsection]{Theorem}
\newtheorem{proposition}[subsection]{Proposition}
\newtheorem{lemma}[subsection]{Lemma}

\theoremstyle{definition}
\newtheorem{definition}[subsection]{Definition}
\newtheorem{example}[subsection]{Example}
\newtheorem{exercise}[subsection]{Exercise}
\newtheorem{situation}[subsection]{Situation}

\theoremstyle{remark}
\newtheorem{remark}[subsection]{Remark}
\newtheorem{remarks}[subsection]{Remarks}

\numberwithin{equation}{subsection}

% Macros
%
\def\lim{\mathop{\mathrm{lim}}\nolimits}
\def\colim{\mathop{\mathrm{colim}}\nolimits}
\def\Spec{\mathop{\mathrm{Spec}}}
\def\Hom{\mathop{\mathrm{Hom}}\nolimits}
\def\Ext{\mathop{\mathrm{Ext}}\nolimits}
\def\SheafHom{\mathop{\mathcal{H}\!\mathit{om}}\nolimits}
\def\SheafExt{\mathop{\mathcal{E}\!\mathit{xt}}\nolimits}
\def\Sch{\mathit{Sch}}
\def\Mor{\mathop{\mathrm{Mor}}\nolimits}
\def\Ob{\mathop{\mathrm{Ob}}\nolimits}
\def\Sh{\mathop{\mathit{Sh}}\nolimits}
\def\NL{\mathop{N\!L}\nolimits}
\def\CH{\mathop{\mathrm{CH}}\nolimits}
\def\proetale{{pro\text{-}\acute{e}tale}}
\def\etale{{\acute{e}tale}}
\def\QCoh{\mathit{QCoh}}
\def\Ker{\mathop{\mathrm{Ker}}}
\def\Im{\mathop{\mathrm{Im}}}
\def\Coker{\mathop{\mathrm{Coker}}}
\def\Coim{\mathop{\mathrm{Coim}}}

% Boxtimes
%
\DeclareMathSymbol{\boxtimes}{\mathbin}{AMSa}{"02}

%
% Macros for moduli stacks/spaces
%
\def\QCohstack{\mathcal{QC}\!\mathit{oh}}
\def\Cohstack{\mathcal{C}\!\mathit{oh}}
\def\Spacesstack{\mathcal{S}\!\mathit{paces}}
\def\Quotfunctor{\mathrm{Quot}}
\def\Hilbfunctor{\mathrm{Hilb}}
\def\Curvesstack{\mathcal{C}\!\mathit{urves}}
\def\Polarizedstack{\mathcal{P}\!\mathit{olarized}}
\def\Complexesstack{\mathcal{C}\!\mathit{omplexes}}
% \Pic is the operator that assigns to X its picard group, usage \Pic(X)
% \Picardstack_{X/B} denotes the Picard stack of X over B
% \Picardfunctor_{X/B} denotes the Picard functor of X over B
\def\Pic{\mathop{\mathrm{Pic}}\nolimits}
\def\Picardstack{\mathcal{P}\!\mathit{ic}}
\def\Picardfunctor{\mathrm{Pic}}
\def\Deformationcategory{\mathcal{D}\!\mathit{ef}}
\usepackage{cleveref}
\crefname{originalsection}{original section}{original sections}
\Crefname{originalsection}{Original section}{Original sections}
\makeatletter
\newlabel{section-relative-fields}{{0H4E}{source}{}{sourcepointer-1}{}}
\newlabel{section-relative-fields@cref}{{[originalsection][0][]0H4E}{[1][0][]source}{}{}{}}
\makeatother
\crefname{originallemma}{original lemma}{original lemmas}
\Crefname{originallemma}{Original lemma}{Original lemmas}
\makeatletter
\newlabel{properties-lemma-locally-Noetherian-specialization-dvr}{{054F}{source}{}{sourcepointer-2}{}}
\newlabel{properties-lemma-locally-Noetherian-specialization-dvr@cref}{{[originallemma][0][]054F}{[1][0][]source}{}{}{}}
\makeatother
\crefname{originalsection}{original section}{original sections}
\Crefname{originalsection}{Original section}{Original sections}
\makeatletter
\newlabel{more-morphisms-section-connected-components}{{055K}{source}{}{sourcepointer-3}{}}
\newlabel{more-morphisms-section-connected-components@cref}{{[originalsection][0][]055K}{[1][0][]source}{}{}{}}
\makeatother
\crefname{originallemma}{original lemma}{original lemmas}
\Crefname{originallemma}{Original lemma}{Original lemmas}
\makeatletter
\newlabel{more-morphisms-lemma-connected-along-section-open}{{055R}{source}{}{sourcepointer-4}{}}
\newlabel{more-morphisms-lemma-connected-along-section-open@cref}{{[originallemma][0][]055R}{[1][0][]source}{}{}{}}
\makeatother
\title{Support of a relative cycle along a smooth morphism of relative dimension r is open and closed}
\author{The Stacks Project Authors}
\date{stacks-project:relative-cycles.tex:local-snapshot-20261012; GNU FDL 1.2 or later}
\begin{document}
\maketitle
\setcounter{subsection}{1}
\noindent\textit{Editorial note.} All delivered mathematical text is the author's own, copied from the pinned Stacks Project chapter file relative-cycles.tex; no statement or proof step was rewritten, reordered or removed. The counted span is the author environment and its own proof environment. References inside the span are handed to the assembler as navigation stubs that carry the canonical tag of the cited tag only; whether each of them is load-bearing is recorded on the stub itself.\par\medskip
\noindent\textit{Allowed standard prerequisites (editorial).}\begin{itemize}
\item Definition of base change of cycles over fields, under which the coefficients of components of a reduced base change are read off from the original cycle (relative-cycles chapter, Section section-relative-fields, canonical tag 0H4E). cited at proof lines 1274-1275. Yields the invariance of the coefficient n\_x under base change: the smooth component Z(x) pulls back to a reduced scheme and Z(x') is an irreducible component of it.
\item Specializations in a locally Noetherian scheme can be realized by a morphism from the spectrum of a discrete valuation ring (Properties chapter, Lemma properties-lemma-locally-Noetherian-specialization-dvr, canonical tag 054F). cited at proof lines 1282-1285. Provides the morphism S' to X mapping the generic point to x' and the closed point to x, turning the specialization check into a valuation-ring computation.
\item Connected components of fibres along a section (More on Morphisms, Section more-morphisms-section-connected-components, canonical tag 055K, and especially Lemma more-morphisms-lemma-connected-along-section-open, canonical tag 055R). cited at proof lines 1295-1298. After replacing X by an open subscheme the fibres of X to S are connected, using the section sigma of the base-changed family.
\end{itemize}\medskip
\ifdefined\hypertarget\hypertarget{sourcepointer-1}{}\fi
\noindent\textit{Source reference.} Original section 0H4E: sources/stacks-chow/tags.txt, source lines 21349--21349. Cited at proof lines 1274-1275: the coefficient invariance under base change rests on the definition of base change given there.\par\medskip
\ifdefined\hypertarget\hypertarget{sourcepointer-2}{}\fi
\noindent\textit{Source reference.} Original lemma 054F: sources/stacks-chow/tags.txt, source lines 6638--6638. Cited at proof lines 1282-1285: realizes the specialization x' leading to x by a valuation ring.\par\medskip
\ifdefined\hypertarget\hypertarget{sourcepointer-3}{}\fi
\noindent\textit{Source reference.} Original section 055K: sources/stacks-chow/tags.txt, source lines 6678--6678. Cited at proof lines 1295-1296: the discussion used to shrink X so its fibres are connected.\par\medskip
\ifdefined\hypertarget\hypertarget{sourcepointer-4}{}\fi
\noindent\textit{Source reference.} Original lemma 055R: sources/stacks-chow/tags.txt, source lines 6684--6684. Cited at proof lines 1297-1298: provides the open subscheme with connected fibres in the DVR case.\par\medskip
\begin{lemma}
\label{lemma-relative-cycle-smooth}
Let $f : X \to S$ be a morphism of schemes. Let $r \geq 0$. Assume $S$
locally Noetherian and $f$ smooth of relative dimension $r$. Let
$\alpha \in z(X/S, r)$. Then the support of $\alpha$ is open and closed in $X$
(see proof for a more precise result).
\end{lemma}

\begin{proof}
Let $x \in X$ with image $s \in S$. Since $f$ is smooth, there is a
unique irreducible component $Z(x)$ of $X_s$ which contains $x$.
Then $\dim(Z(x)) = r$.
Let $n_x$ be the coefficient of $Z(x)$ in the cycle $\alpha_s$.
We will show the function $x \mapsto n_x$ is locally constant on $X$.

\medskip\noindent
Let $g : S' \to S$ be a morphism of locally Noetherian schemes.
Let $X'$ be the base change of $X$ and let $\alpha' = g^*\alpha$
be the base change of $\alpha$. Let $x' \in X'$ map to $s' \in S'$,
$x \in X$, and $s \in S$. We claim $n_{x'} = n_x$. Namely, since $Z(x)$
is smooth over $\kappa(s)$ we see that
$Z(x) \times_{\Spec(\kappa(s))} \Spec(\kappa(s'))$
is reduced. Since $Z(x')$ is an irreducible component
of this scheme, we see that the coefficient $n_{x'}$ of $Z(x')$
in $\alpha'_{s'}$ is the same as the coefficient $n_x$ of $Z(x)$
in $\alpha_s$ by the definition of base change in
Section \ref{section-relative-fields}
thereby proving the claim.

\medskip\noindent
Since $X$ is locally Noetherian, to show that $x \mapsto n_x$
is locally constant, it suffices to show: if $x' \leadsto x$
is a specialization in $X$, then $n_{x'} = n_x$.
Choose a morphism $S' \to X$ where $S'$ is the spectrum of a discrete
valuation ring mapping the generic point $\eta$ to $x'$ and the closed point
$0$ to $x$. See Properties, Lemma
\ref{properties-lemma-locally-Noetherian-specialization-dvr}.
Then the base change $X' \to S'$ of $f$ by $S' \to S$
has a section $\sigma : S' \to X'$ such that $\sigma(\eta) \leadsto \sigma(0)$
is a specialization of points of $X'$ mapping to $x' \leadsto x$ in $X$.
Thus we reduce to the claim in the next paragraph.

\medskip\noindent
Let $S$ be the spectrum of a discrete valuation ring with
generic point $\eta$ and closed point $0$ and we have a section
$\sigma : S  \to X$. Claim: $n_{\sigma(\eta)} = n_{\sigma(0)}$.
By the discussion in More on Morphisms, Section
\ref{more-morphisms-section-connected-components} and especially
More on Morphisms, Lemma \ref{more-morphisms-lemma-connected-along-section-open}
after replacing $X$ by an open subscheme, we may assume the
fibres of $X \to S$ are connected. Since these fibres
are smooth, they are irreducible. Then we see that $\alpha_\eta = n[X_\eta]$
with $n = n_{\sigma(\eta)}$ and the relation $sp_{X/S}(\alpha_\eta) = \alpha_0$
implies $\alpha_0 = n[X_0]$, i.e., $n_{\sigma(0)} = n$ as desired.
\end{proof}

\end{document}
